\section{迭代：把解写成一列近似的极限}
\label{sec:17-Synthesis/02-Recurring-Ideas/03}

一个价值迭代的实验。状态是表格形式时，\(\gamma=0.99\)，几千轮之后价值函数稳稳停住，误差按 \(0.99^{k}\) 掉下去，与理论完全对得上。把表格换成一个线性函数逼近器，其余一行不改，同一个 \(\gamma\)，同一份转移数据——几万步之后价值估计爬到 \(10^{7}\)，然后是 NaN。

先怀疑学习率，调小，只是发散得慢一点。再怀疑梯度实现，查了一遍，没错。真正的原因不在代码里：表格情形下 Bellman 算子在无穷范数下是 \(\gamma\) 压缩，而``先做 Bellman 更新、再投影回可表示的函数类''这个复合算子，压缩性只在某个特定的加权范数下成立，权重由行为策略的访问分布给出。数据不是按那个分布采的，复合算子就可能不是压缩，迭代于是没有任何理由收敛。

这件事的教训不是关于强化学习的。它关于一个更普遍的动作，以及那个动作真正依赖的东西。

\subsection{一个动作}

第三个反复出现的数学动作是：\emph{解不出来的方程，改写成``求某个映射的不动点''，然后从任意一点开始反复作用这个映射。}把 \(g(x)=0\) 写成 \(x = T(x)\)，构造序列 \(x_{k+1}=T(x_k)\)，把解定义成这列近似的极限。它把``求解''这件事从代数问题换成了收敛问题，而收敛问题有一整套现成的判据。

这个改写只有在两件事都成立时才有意义：极限存在，而且极限还在原来的空间里。后一条不是废话，\hyperref[sec:12-Real-Analysis-Support/03-Topology-Basics/03]{``完备性让逼近留在空间内''}给出了它失效的形态——在有理数里迭代逼近 \(\sqrt2\)，每一项都合法，极限却不在空间中，整个论证白做。所有迭代算法的收敛论证最终都要落在某个完备空间上，这是完备性在这本书里被反复要求的原因。\hyperref[sec:12-Real-Analysis-Support/01-Real-Numbers-and-Completeness/02]{``Cauchy 判据无需预知极限''}给的则是使用它的方式：不必先知道答案是什么，只看序列自己后面的项挤不挤在一起。

\begin{center}
\begin{tikzpicture}[x=1cm,y=1cm,font=\small,>=stealth]
  % ---- 左：压缩映射下误差球等比缩小 ----
  \draw[black!25] (2.2,2.2) circle (1.9);
  \draw[black!30] (2.2,2.2) circle (1.14);
  \draw[black!40] (2.2,2.2) circle (0.684);
  \draw[black!55] (2.2,2.2) circle (0.41);
  \fill[red!75!black] (2.2,2.2) circle (0.07);
  \node[font=\footnotesize,red!70!black,anchor=south] at (2.15,2.30) {\(x^{\star}\)};
  \draw[blue!60!black,thick,->] (3.99,2.85) -- (1.33,2.93);
  \draw[blue!60!black,thick,->] (1.33,2.93) -- (2.08,1.53);
  \draw[blue!60!black,thick,->] (2.08,1.53) -- (2.59,2.34);
  \foreach \p in {(3.99,2.85),(1.33,2.93),(2.08,1.53),(2.59,2.34)} \fill[blue!60!black] \p circle (0.055);
  \node[font=\footnotesize,black!70,anchor=west] at (3.9,3.15) {\(x_0\)};
  \node[font=\footnotesize,black!70,anchor=south] at (2.2,4.3) {半径每步乘 \(\gamma<1\)};
  \node[font=\footnotesize,black!70] at (2.2,-0.15) {压缩：存在、唯一、几何收敛};

  % ---- 右：不动点与真正想要的解不重合 ----
  \begin{scope}[shift={(6.4,0)}]
    \draw[black!55,thick] (0.2,0.8) -- (5.0,2.2);
    \node[font=\footnotesize,black!70,anchor=north west] at (3.9,0.95) {可表示的函数类};
    \fill[red!75!black] (2.4,3.6) circle (0.07);
    \node[font=\footnotesize,red!70!black,anchor=south] at (2.4,3.72) {真解};
    \draw[black!45,dashed] (2.4,3.6) -- (2.98,1.61);
    \fill[green!45!black] (2.98,1.61) circle (0.07);
    \draw[green!35!black,thin,dotted] (2.95,1.53) -- (2.75,1.22);
    \node[font=\footnotesize,green!35!black,anchor=north] at (2.65,1.18) {最佳逼近};
    \fill[blue!60!black] (4.3,2.0) circle (0.07);
    \node[font=\footnotesize,blue!60!black,anchor=north west] at (4.45,1.78) {迭代的不动点};
    \draw[orange!85!black,very thick,|-|] (2.98,1.61) -- (4.3,2.0);
    \node[font=\footnotesize,orange!70!black,anchor=south] at (3.64,1.86) {差距};
    \node[font=\footnotesize,black!70] at (2.6,-0.15) {收敛不等于收敛到你要的东西};
  \end{scope}
\end{tikzpicture}

\emph{左边是这个动作最好的情形：误差球按固定比例缩小，起点在哪里都不重要。右边是它最容易被忽略的代价——迭代收敛到的那个点由算子决定，而不是由``离真解最近''决定，两者可以差得很远。}
\end{center}

\subsection{压缩给出的三样东西}

若 \(T\) 在某个完备度量空间上满足 \(d(T(x),T(y)) \le \gamma\, d(x,y)\) 且 \(\gamma<1\)，则不动点存在、唯一，且从任意初值出发都以 \(\gamma^{k}\) 的速率收敛。这三样是打包给的，不能只要一样。存在性来自 Cauchy 列加完备；唯一性来自两个不动点之间的距离必须小于自己；速率直接从定义读出。

还有一样常被忘记的赠品：可用的停机准则。由三角不等式，
\[
  d(x_k, x^{\star}) \;\le\; \frac{\gamma}{1-\gamma}\, d(x_k, x_{k-1}),
\]
右端全是能算的量。这就是为什么迭代算法可以拿``相邻两步的变化''当误差指标——它不是启发式，是压缩性的推论。而 \(\gamma\) 越接近 \(1\)，这个因子越大，同样的``看上去不动了''对应的真实误差越大。折扣因子 \(0.99\) 与 \(0.9\) 在这个式子里差十倍，而经常有人只把它当成一个``偏好未来的程度''。

\hyperref[sec:08-Numerical-Analysis/02-Root-Finding/02]{``不动点与 Newton 法''}把这套判据用在最基本的场合，并说明 Newton 法为什么快得多——它的迭代映射在解处导数为零，压缩模数局部趋于零，收敛从几何变成二次。\hyperref[sec:15-Differential-Equations-and-Dynamical-Systems/01-ODE-Theory/03]{``存在唯一性为什么需要 Lipschitz''}给出同一套判据的另一处使用：微分方程的解被写成一个积分算子的不动点，Lipschitz 条件正是为了让那个算子在短时间区间上成为压缩。存在唯一性定理与迭代算法在这里是同一件事的两种读法。

\subsection{同一个动作的其它名字}

值迭代是这套理论最干净的应用之一，干净到 \(\gamma\) 直接就是压缩模数。\hyperref[sec:16-Control-and-Reinforcement-Learning/02-Reinforcement-Learning/02]{``Bellman 方程与值迭代、策略迭代''}说明 Bellman 最优算子在无穷范数下是 \(\gamma\) 压缩，于是最优价值函数的存在唯一、以及从任意初始估计出发的收敛，全部是上一段那条定理的直接推论——不需要任何强化学习特有的论证。

幂法把``求主特征向量''写成迭代：反复乘矩阵再归一化，收敛速率由前两个特征值之比决定，那个比值就是这里的 \(\gamma\)。\hyperref[sec:08-Numerical-Analysis/05-Numerical-Linear-Algebra/05]{``从幂迭代到 QR 特征值算法''}顺着这条线走到实用算法。求解大型线性方程组时同样如此：Jacobi 与 Gauss--Seidel 把 \(Ax=b\) 拆成 \(x = Mx + c\)，收敛与否取决于 \(M\) 的谱半径小于一，与压缩条件是同一句话的矩阵版本；\hyperref[sec:08-Numerical-Analysis/05-Numerical-Linear-Algebra/04]{``大型稀疏系统与 Krylov 迭代''}讲的是当直接分解在规模上不可行时，为什么迭代反而成为唯一选择。

Markov 链的平稳分布是转移算子的不动点，\hyperref[sec:06-Random-Processes/03-Markov-Chains/04]{``平稳分布、可逆性与遍历收敛''}给出它存在唯一并被逼近的条件；\hyperref[sec:13-Machine-Learning/08-Sampling-and-Inference/02]{``MCMC 在无法独立采样时构造平稳链''}把这件事倒过来用——不是给定链求平稳分布，而是给定想要的分布去造一条以它为不动点的链，然后跑到收敛。混合时间在这里扮演的正是 \(\gamma\) 的角色，而它通常估不出来，这是 MCMC 诊断永远不能给出证明的根本原因。

训练本身也是这个动作。梯度下降是映射 \(x\mapsto x-\eta\nabla f(x)\) 的迭代，在强凸加光滑的条件下它确实是压缩，压缩模数由条件数给出；上一篇讲的换基在这里派上用场——那个模数正是特征值的极端比。深度学习的训练没有这样的保证，但形式完全一致，所以关于步长、关于停机、关于``看上去不动了''的所有直觉都可以搬过来，只是必须记住它们此时是经验规律而非定理。

\subsection{三种不同的保证}

压缩不是唯一的收敛理由，另外两条在这本书里出现得同样频繁，给出的保证却弱得多。

第一条是单调。EM 算法每一步都不降低似然，\hyperref[sec:13-Machine-Learning/06-Probabilistic-Models/04]{``EM 用下界交替完成隐变量学习''}把这个性质的来源讲清：E 步构造一个在当前参数处相切的下界，M 步最大化这个下界，于是似然至少不降。有界加单调给出收敛，但只给出收敛——不给唯一性，不给速率，也不保证收敛到全局最优。混合模型跑十次得到十个不同的解，不是实现有问题，是这条保证本来就只承诺这么多。

第二条是随机逼近。每一步只看到一个带噪声的样本，\hyperref[sec:16-Control-and-Reinforcement-Learning/02-Reinforcement-Learning/03]{``Monte Carlo 与随机逼近把回报变成增量更新''}给出 Robbins--Monro 的两个条件：
\[
  \sum_{t} \alpha_t = \infty, \qquad \sum_{t} \alpha_t^{2} < \infty .
\]
这两个条件各挡一件事，值得分开读。第一条挡的是``走不到''——步长和收敛的话，无论初值离目标多远，总步进量有限，可能永远到不了。第二条挡的是``停不下''——步长平方和发散的话，噪声的累积方差不衰减，迭代会一直在解附近抖动。合起来就是``走得够远，但抖得越来越小''。\hyperref[sec:09-Convex-Optimization/08-Nonconvex-and-Stochastic-Optimization/04]{``SGD 在期望中下降但噪声会留下平台''}说明固定步长违反第二条时会发生什么：损失下降到某个水平就停在那里，那个平台的高度正比于步长与梯度噪声方差。深度学习里的学习率衰减就是在事后补上第二条，只不过实践中用的衰减速度通常比理论要求快得多，这属于经验做法。

第三条是压缩，也就是前面那一套。三者的强弱是清楚的：压缩给存在唯一加速率，单调只给收敛，随机逼近在噪声下给概率一收敛但速率对步长极其敏感。写代码时用的是同一个 \texttt{while} 循环，能承诺的东西差三个量级。

\subsection{收敛到不动点，不等于收敛到你想要的东西}

这是这个动作最容易失守的地方，也是开头那个故事的真正内容。

把价值函数限制在一个线性函数类里，迭代收敛到的是``先做 Bellman 更新再投影''这个复合算子的不动点。\hyperref[sec:16-Control-and-Reinforcement-Learning/02-Reinforcement-Learning/05]{``线性价值函数逼近求解投影 Bellman 方程''}说明它一般\emph{不是}真价值函数在该函数类中的最佳逼近，两者的差距有一个带 \(1/(1-\gamma)\) 因子的界——折扣接近一时这个界近乎无用。上面那张图的右半边画的就是这件事：绿点是你以为会得到的，蓝点是实际得到的。

而且压缩性本身依赖于范数的选择。表格情形下 Bellman 算子在无穷范数下压缩；带投影之后，压缩只在以行为策略访问分布加权的范数下成立。数据来自另一个策略时，投影是在错的权重下做的，复合算子的模数可以超过一，迭代不再收敛。\hyperref[sec:16-Control-and-Reinforcement-Learning/02-Reinforcement-Learning/06]{``Deadly Triad 让自举误差在分布外放大''}把函数逼近、自举、离策略这三件事凑齐时的发散机制讲透了。三件事里任意去掉一件通常就没事，这也说明问题不在其中某一件，而在它们对压缩性所依赖的那个范数的联合破坏。

同样的事在别处换个说法出现。EM 收敛到似然的局部极大，初值决定去哪一个。MCMC 收敛到目标分布，但从有限长的链上根本判断不出是否已经收敛。Newton 法收敛到某个根，未必是你要的那个根。这些都不是算法缺陷，是这个动作本身的性质：迭代只承诺把你送到算子的不动点，``那个不动点是不是解''是另一个问题，必须单独回答。

\subsection{边界}

值得记住的判断可以压成一句：\emph{只要看到一个循环里反复作用同一个映射，就该问三件事——它在什么范数下压缩、模数是多少、以及它的不动点是不是你要的解。}三个问题各自对应上面三节。多数时候前两个有答案而第三个没有，这时应当明说保证到哪一层为止，而不是把``收敛了''直接汇报成``解出来了''。

这一篇与前两篇共享同一个形状：都是先把难问题换成便宜的形式，再交代这次换算的有效范围。下一篇换一个方向——不去构造解，而是去证明某些东西根本不可能出现。那类论证不需要迭代，也不需要好坐标，只需要找到一个在变化中不动的量。
